\section*{الفصل \ref{ch:g4:animal-reproduction} --- كيف تتكاثر الحيوانات}

\begin{solution}{exo:g4:animal-reproduction:1}
التحام \omterm{def:g4:animal-reproduction:sexual}{نطفة} واحدة \omterm{def:g4:animal-reproduction:sexual}{ببويضة} واحدة، وهو الذي يبدأ حيوانًا صغيرًا
جديدًا. \omterm{def:g4:animal-reproduction:sexual}{والبويضة} تأتي من \omterm{def:g4:animal-reproduction:sexual}{الأنثى}، \omterm{def:g4:animal-reproduction:sexual}{والنطفة} من \omterm{def:g4:animal-reproduction:sexual}{الذكر}.
\end{solution}

\begin{solution}{exo:g4:animal-reproduction:2}
\omterm{def:g4:animal-reproduction:ovivivi}{البيوض}: \omterm{def:g4:animal-reproduction:sexual}{الأنثى} تضع بيضًا ويتم الصغير نموه في \omterm{def:g2:animals-grow-up:born}{البيضة} خارج جسمها ---
كالدجاجة والضفدع (وكذلك السلمون المرقط والسحالي والحلزونات
\omterm{prop:g3:sorting-living-things:groups}{والحشرات}). \omterm{def:g4:animal-reproduction:ovivivi}{والولود}: الصغير ينمو داخل الأم مغذًّى منها ثم \omterm{def:g2:animals-grow-up:born}{يولد حيًّا}
--- كالفرس والحوت (\omterm{prop:g3:sorting-living-things:groups}{والثدييات} كلها تقريبًا).
\end{solution}

\begin{solution}{exo:g4:animal-reproduction:3}
غنّت الذكور لتنادي الإناث إلى البركة؛ ثم أطلق كل زوج خلاياه في
الماء حيث لاقت النطف \omterm{def:g4:animal-reproduction:sexual}{البويضات} --- وهو \omterm{def:g4:animal-reproduction:sexual}{الإخصاب}؛ ونمت \omterm{def:g4:animal-reproduction:ovivivi}{البيوض}
المخصَّبة في هلامها فصارت \omterm{ex:g2:animals-grow-up:frog}{شراغيف} حزيران.
\end{solution}

\begin{solution}{exo:g4:animal-reproduction:4}
\omterm{def:g4:animal-reproduction:sexual}{إخصاب} السلمون المرقط يحدث في الماء خارج \omterm{def:g4:animal-reproduction:sexual}{الأنثى} --- إذ تُلقى
\omterm{def:g4:animal-reproduction:sexual}{البويضات} والنطف معًا. \omterm{def:g4:animal-reproduction:sexual}{وإخصاب} الدجاجة يحدث داخل جسمها في أثناء
\omterm{prop:g4:animal-reproduction:where}{التزاوج}. فعلى اليابسة تجفّ الخلايا العارية؛ أما في الماء فتستطيع
أن تلتقي بأمان في الخارج.
\end{solution}

\begin{solution}{exo:g4:animal-reproduction:5}
الكتكوت ينمو في \omterm{def:g2:animals-grow-up:born}{بيضة} مقشورة خارج الدجاجة، عائشًا على المحّ المحزوم
معه. والمهر ينمو أحد عشر شهرًا داخل الفرس، مغذًّى بجسمها باستمرار،
ثم \omterm{def:g2:animals-grow-up:born}{يولد حيًّا}.
\end{solution}

\begin{solution}{exo:g4:animal-reproduction:6}
كلما أنتج النوع صغارًا أكثر قلّت العناية بكل واحد؛ وكلما أنتج أقلّ
زادت العناية. والطرفان: ملايين \omterm{def:g4:animal-reproduction:ovivivi}{بيوض} سمكة القدّ غير المعتنى بها،
وعجل الفيلة الواحد المحروس عقدًا كاملًا.
\end{solution}

\begin{solution}{exo:g4:animal-reproduction:7}
لأن النجاح لا يقتضي إلا أن يبقى في المتوسط من الصغار ما يكفي ليحلّ
محل الوالدين. فمن بين ملايين \omterm{def:g4:animal-reproduction:ovivivi}{البيوض} تفلت دائمًا قلة من أخطار البحر
--- والعدد الهائل هو \emph{نفسه} الحماية.
\end{solution}

\begin{solution}{exo:g4:animal-reproduction:8}
صغار قليلة (فخمسة قليل بالقياس إلى الآلاف)، نمت داخل الأم،
وتُرضَع وتُحرَس: فالقطة تقع بوضوح في جانب القليل المعتنى به ---
\omterm{prop:g3:sorting-living-things:groups}{كالثدييات} عمومًا.
\end{solution}

\begin{solution}{exo:g4:animal-reproduction:9}
على أنه ينال شيئًا من \emph{كل} واحد من والديه --- فالصغير ليس
نسخة من أمه وحدها. وكلا الوالدين يسهم فيما يصير إليه المهر.
\end{solution}

\begin{solution}{exo:g4:animal-reproduction:10}
حضنة صغيرة --- بضع عشرات أو مئات من \omterm{def:g4:animal-reproduction:ovivivi}{البيوض}، لا ملايين. فالعناية
والعدد يتبادلان: فالأب الذي يحرس ويروّح ويدافع يستطيع أن يعبر
بحضنة صغيرة أخطارًا كانت ستقتضي ملايين \omterm{def:g4:animal-reproduction:ovivivi}{البيوض} غير المحروسة.
\end{solution}

\begin{solution}{exo:g4:animal-reproduction:11}
\omterm{def:g4:animal-reproduction:ovivivi}{بيوض} --- وغير معتنًى ببيضها: فالحضنة المدفونة تُترك لدفء الرمل،
وتنطلق الفواقس إلى البحر وحدها (وعددها الكبير هو خطتها). وعليها أن
تضع على اليابسة لأن \omterm{def:g2:animals-grow-up:born}{بيض} \omterm{prop:g4:classifying-animals:classes}{الزواحف} مقشور ويتنفس الهواء --- فلو وُضع
في البحر لغرق. فصفات طائفتها، لا عاداتها، هي التي تضع القاعدة.
\end{solution}
